\section{L-symbols of a transitive action from stabilizer data}
\label{sec:threecocycle_stabilizer}

For a transitive action of $G$ on $\mathsf X$ with a base point $x_0$ and
stabilizer $H$, chosen representatives of the cosets turn every two-cocycle
on $H$ into an $L$-symbol, and every $L$-symbol compatible with a
three-cocycle that is trivial on $H$ arises in this way up to a gauge. The
gauge classes of $L$-symbols compatible with a fixed three-cochain form a
torsor under $H^2(H,\C^\times)$, following Garre-Rubio, Lootens and
Moln\'ar~\cite{GarreRubio2023MPOSym}. An $L$-symbol induced from $H$ can be
trivialized by a gauge exactly when the class of its cocycle extends to $G$,
following~\cite{FrancoRubio2025MPUGauging}; for finite $G$ the criterion
holds equally with unit-circle coefficients.

\begin{lemma}[Stabilizer transition elements]
    \label{thm:mpug_stabilizer_transition}
    \lean{StabilizerRepresentatives.transitionElement_mul,
        StabilizerRepresentatives,
        StabilizerRepresentatives.ofPretransitive,
        StabilizerRepresentatives.transitionElement,
        StabilizerRepresentatives.coe_transitionElement,
        StabilizerRepresentatives.transitionElement_stabilizer}
    \leanok
    Suppose that $G$ acts transitively on $\mathsf X$, and fix
    $x_0\in\mathsf X$. There are representatives $k_x\in G$ satisfying
    $k_xx_0=x$ and $k_{x_0}=e$. With
    $H=\operatorname{Stab}_G(x_0)$, define
    \begin{align}
        t_k(g,x)
         & =k_{g x}^{-1}gk_x.
        \label{eq:mpug_stabilizer_transition}
    \end{align}
    Then $t_k(g,x)\in H$ and
    \begin{align}
        t_k(g_1g_2,x)
         & =t_k(g_1,g_2x)t_k(g_2,x).
        \label{eq:mpug_stabilizer_transition_product}
    \end{align}
    Consequently, the source factors are
    \begin{align}
        h_2
         & =t_k(g_2,x)=k_{g_2x}^{-1}g_2k_x,
        \label{eq:mpug_stabilizer_h2}                  \\
        h_1
         & =t_k(g_1,g_2x)=k_{g_1g_2x}^{-1}g_1k_{g_2x}.
        \label{eq:mpug_stabilizer_h1}
    \end{align}
    As recorded in \cite{gap:fbc25_stabilizer_transition_index_typo}, the last
    representative in $h_1$ corrects the printed $k_{g_1x}$. No finiteness or
    normality assumption is required.
    % Local correction:
    % docs/paper-gaps/fbc25_stabilizer_transition_index_typo.tex.
    % Source anchor: arXiv:2502.20257, lemma:h1h2 and eq:h1h2,
    % lines 7028--7042.
\end{lemma}
\begin{proof}\leanok
    Transitivity gives one element carrying $x_0$ to each $x$, with the
    identity chosen at $x=x_0$. Acting with (\ref{eq:mpug_stabilizer_transition})
    on $x_0$ gives
    $k_{gx}^{-1}g(k_xx_0)=k_{gx}^{-1}(gx)=x_0$.
    For (\ref{eq:mpug_stabilizer_transition_product}), substituting
    (\ref{eq:mpug_stabilizer_transition}) and cancelling
    $k_{g_2x}k_{g_2x}^{-1}$ yields the product formula. The two formulas for
    $h_1$ and $h_2$ are the corresponding specializations.
\end{proof}

\begin{lemma}[L-symbol induced by a stabilizer cocycle]
    \label{thm:mpug_stabilizer_cocycle_l_symbol}
    \lean{StabilizerRepresentatives.inducedLSymbol_isCompatible,
        StabilizerRepresentatives.inducedLSymbol,
        StabilizerRepresentatives.inducedLSymbol_base,
        StabilizerRepresentatives.inducedActionGauge,
        StabilizerRepresentatives.inducedLSymbol_eq_gauge,
        StabilizerRepresentatives.isMulCoboundary₂_div_iff_exists_actionGauge_inducedLSymbol_eq}
    \leanok
    \uses{def:mpug_l_symbols, def:mpug_l_gauge, def:is_cocycle, def:cohomologous_to, thm:mpug_stabilizer_transition}
    Let $H=\operatorname{Stab}_G(x_0)$ and let
    $\psi\colon H\times H\to\C^\times$ be a scalar $2$-cocycle. Define the
    $\omega=1$ scalar specialization of \cite[Eq.~(20)]{GarreRubio2023MPOSym}
    by
    \begin{align}
        L^x_{g,h}
         & =\psi\bigl(t_k(g,hx),t_k(h,x)\bigr).
        \label{eq:mpug_induced_stabilizer_l_symbol}
    \end{align}
    The $L$-symbol (\ref{eq:mpug_induced_stabilizer_l_symbol}) satisfies
    (\ref{eq:mpug_l_compatibility}) with $\omega=1$, and its base-point
    restriction recovers $\psi$:
    \begin{align}
        L^{x_0}_{a,b} & =\psi(a,b)
        \label{eq:mpug_induced_stabilizer_base}
    \end{align}
    for all $a,b\in H$.

    More precisely, two stabilizer cocycles are cohomologous exactly when
    their induced $L$-symbols differ by a gauge one-cochain and the trivial
    two-cochain:
    \begin{align}
        \psi_1\sim\psi_2
         & \quad\Longleftrightarrow\quad
        \exists\gamma,\quad
        L[\psi_1]=(1,\gamma)\mathbin{\triangleright}L[\psi_2].
        \label{eq:mpug_induced_cohomology_iff}
    \end{align}
    If the cohomology relation is given by
    $\varphi\colon H\to\C^\times$, with
    \begin{align}
        \psi_1(a,b)
         & =\varphi(a)\varphi(b)\varphi(ab)^{-1}\psi_2(a,b),
        \label{eq:mpug_stabilizer_cocycle_gauge}
    \end{align}
    then the forward gauge one-cochain is
    \begin{align}
        \gamma_{g,x} & =\varphi\bigl(t_k(g,x)\bigr)^{-1}.
        \label{eq:mpug_induced_action_gauge}
    \end{align}
    Thus the induced action-gauge class depends only on the cohomology class
    of $\psi$.

    This is the $L$-symbol form of the reduction in
    \cite[lines~7028--7042]{FrancoRubio2025MPUGauging}.
    % Source anchor: arXiv:2203.12563, Eq. (20).
\end{lemma}
\begin{proof}\leanok
    \uses{def:mpug_l_gauge, thm:mpug_stabilizer_transition}
    With
    $a=t_k(g,hkx)$, $b=t_k(h,kx)$, and $c=t_k(k,x)$,
    the product law (\ref{eq:mpug_stabilizer_transition_product}) gives
    $t_k(gh,kx)=ab$ and $t_k(hk,x)=bc$. Compatibility is therefore
    \begin{align}
        \psi(a,b)\psi(ab,c)
         & =\psi(a,bc)\psi(b,c),
    \end{align}
    which is the cocycle equation. If $p,q\in H$, then
    $t_k(p,x_0)=p$ and $t_k(q,x_0)=q$ because $k_{x_0}=e$, proving
    (\ref{eq:mpug_induced_stabilizer_base}). For the forward gauge
    calculation, setting $a=t_k(g,hx)$ and $b=t_k(h,x)$ gives
    $t_k(gh,x)=ab$. Substituting (\ref{eq:mpug_induced_action_gauge}) into
    (\ref{eq:mpug_l_gauge}) gives the multiplier
    \begin{align}
        \frac{\gamma_{gh,x}}{\gamma_{g,hx}\gamma_{h,x}}
         & =\varphi(ab)^{-1}
        \bigl(\varphi(a)^{-1}\varphi(b)^{-1}\bigr)^{-1}
        =\varphi(a)\varphi(b)\varphi(ab)^{-1}.
    \end{align}
    Formula (\ref{eq:mpug_stabilizer_cocycle_gauge}) now identifies the two
    induced $L$-symbols. Conversely, if
    $L[\psi_1]=(1,\gamma)\mathbin{\triangleright}L[\psi_2]$ and
    $\varphi(a)=\gamma_{a,x_0}^{-1}$ for $a\in H$, evaluating the gauge
    equality at $x_0$ and $a,b\in H$ gives
    \begin{align}
        \psi_1(a,b)
         & =\frac{\gamma_{ab,x_0}}
            {\gamma_{a,x_0}\gamma_{b,x_0}}\psi_2(a,b)
        =\varphi(a)\varphi(b)\varphi(ab)^{-1}\psi_2(a,b).
    \end{align}
    Hence $\psi_1\sim\psi_2$, which proves the reverse implication in
    (\ref{eq:mpug_induced_cohomology_iff}).
\end{proof}

\begin{lemma}[Reconstruction of stabilizer L-symbols]
    \label{thm:mpug_stabilizer_reconstruction}
    \lean{StabilizerRepresentatives.reconstructedLSymbol,
        StabilizerRepresentatives.reconstructedLSymbol_isCompatible,
        StabilizerRepresentatives.reconstructedLSymbol_one_isCompatible,
        StabilizerRepresentatives.reconstructedLSymbol_eq_mul,
        StabilizerRepresentatives.reconstructedLSymbol_base,
        StabilizerRepresentatives.gauge_reconstructedLSymbol_isCompatible,
        ScalarThreeCochain.exists_fusionGauge_eq_one_of_isTrivialGaugeClass_comap,
        LSymbol.IsCompatible.mul}
    \leanok
    \uses{def:mpug_l_symbols, def:mpug_l_gauge, def:mpug_three_cocycle, def:mpug_fusion_gauge, def:mpug_trivial_gauge_class, def:is_cocycle, thm:mpug_stabilizer_transition}
    Let $G$ act on $\mathsf X$ with basepoint $x_0$, stabilizer
    $H=\operatorname{Stab}_G(x_0)$, and representatives $k_x$ as in
    Lemma~\ref{thm:mpug_stabilizer_transition}. Let $\omega$ be a scalar
    $3$-cocycle on $G$ with $\omega(a,b,c)=1$ for all $a,b,c\in H$, and let
    $\psi$ be a scalar $2$-cocycle on $H$. With $h_2=t_k(g_2,x)$ and
    $h_1=t_k(g_1,g_2x)$, define
    \begin{align}
        L^x_{g_1,g_2}
         & =\frac{\omega(g_1g_2k_x,h_2^{-1},h_1^{-1})\,\psi(h_1,h_2)}
            {\omega(g_1,g_2k_x,h_2^{-1})\,\omega(g_1,g_2,k_x)}.
        \label{eq:mpug_stabilizer_reconstruction}
    \end{align}
    Then the $L$-symbol (\ref{eq:mpug_stabilizer_reconstruction}) of
    \cite[Eq.~(20)]{GarreRubio2023MPOSym} satisfies
    (\ref{eq:mpug_l_compatibility}), with
    $L^{x_0}_{a,b}=\psi(a,b)$ for $a,b\in H$.

    If only $\omega|_H$ is a coboundary, some scalar $2$-cochain $\beta$ on $G$
    makes $\beta\mathbin{\triangleright}\omega$ equal to one on
    $H\times H\times H$, and applying the inverse gauge
    $(\beta^{-1},1)$ to (\ref{eq:mpug_stabilizer_reconstruction}) built from
    $\beta\mathbin{\triangleright}\omega$ gives an $L$-symbol compatible with
    $\omega$. The hypothesis that $\omega$ is one on the stabilizer fixes the
    gauge in which the source's formula holds; as recorded in
    \cite{gap:glm23_eq20_fusion_gauge}, for a general
    $\omega$ with $\omega|_H=d\beta$ the printed formula can fail.
    No normalization of $\omega$ is required.
    % Local fix: docs/paper-gaps/glm23_eq20_fusion_gauge.tex.
    % Source anchor: arXiv:2203.12563, REsubmission.tex lines 740--761.
\end{lemma}
\begin{proof}\leanok
    \uses{thm:mpug_stabilizer_transition, thm:mpug_stabilizer_cocycle_l_symbol}
    The formula factors as its value at $\psi=1$ times the induced
    $L$-symbol (\ref{eq:mpug_induced_stabilizer_l_symbol}), and pointwise
    products of compatible $L$-symbols are compatible with the product of
    the $3$-cochains. By Lemma~\ref{thm:mpug_stabilizer_cocycle_l_symbol}
    it suffices to treat $\psi=1$. For fixed $x$ and $g,h,k$, with $p=k_x$,
    $a=t_k(k,x)$, $b=t_k(h,kx)$, and $c=t_k(g,hkx)$, we have $k_{kx}=kpa^{-1}$
    and, by (\ref{eq:mpug_stabilizer_transition_product}),
    $t_k(hk,x)=ba$ and $t_k(gh,kx)=cb$. The compatibility equation then
    follows from the $3$-cocycle equation at the four quadruples
    \begin{align}
         & (g,h,k,p),\quad (g,h,kp,a^{-1}),\quad (g,hkp,a^{-1},b^{-1}),\quad
        (ghkp,a^{-1},b^{-1},c^{-1}),
    \end{align}
    together with $\omega(a^{-1},b^{-1},c^{-1})=1$. At $x_0$ and $a,b\in H$
    every argument of $\omega$ in (\ref{eq:mpug_stabilizer_reconstruction})
    lies in $H$, which gives $L^{x_0}_{a,b}=\psi(a,b)$. If
    $\omega|_H=d\beta_0$, extending $\beta_0^{-1}$ by one outside
    $H\times H$ produces a coboundary that cancels $\omega$ on
    $H\times H\times H$. Finally, $(\beta^{-1},1)$ changes
    $\beta\mathbin{\triangleright}\omega$ back to $\omega$.
\end{proof}

\begin{lemma}[Classification and torsor of stabilizer L-symbols]
    \label{thm:mpug_stabilizer_torsor}
    \lean{StabilizerRepresentatives.exists_actionGauge_eq_inducedLSymbol,
        StabilizerRepresentatives.exists_actionGauge_eq_reconstructedLSymbol,
        StabilizerRepresentatives.exists_actionGauge_gauge_eq_reconstructedLSymbol,
        StabilizerRepresentatives.isCompatible_iff_exists_actionGauge_mul_inducedLSymbol,
        StabilizerRepresentatives.actionGaugeEquiv_mul_inducedLSymbol_iff,
        StabilizerRepresentatives.solutionSetoid,
        StabilizerRepresentatives.h2EquivSolutionClasses,
        StabilizerRepresentatives.h2EquivSolutionClasses_H2π,
        StabilizerRepresentatives.h2EquivSolutionClasses_apply,
        StabilizerRepresentatives.solutionAddAction,
        StabilizerRepresentatives.existsUnique_vadd_eq,
        LSymbol.ActionGaugeEquiv,
        LSymbol.actionGaugeEquiv_equivalence,
        LSymbol.IsCompatible.inv,
        LSymbol.IsCompatible.mul_inv_one,
        LSymbol.IsCompatible.gauge_one,
        LSymbol.IsCompatible.isCocycle_restrict,
        LSymbol.gauge_one_mul}
    \leanok
    \uses{thm:mpug_stabilizer_reconstruction, def:mpug_l_gauge, thm:mpug_h2_group_cohomology_comparison, def:cohomologous_to}
    Keep the setting of Lemma~\ref{thm:mpug_stabilizer_reconstruction},
    and call two $L$-symbols action-gauge equivalent when they differ by
    $(1,\gamma)\mathbin{\triangleright}$ for some gauge one-cochain
    $\gamma$.
    \begin{enumerate}
        \item If $R$ is compatible with the trivial $3$-cochain, then
              $R=(1,\gamma)\mathbin{\triangleright}L[\psi]$ for the induced
              $L$-symbol (\ref{eq:mpug_induced_stabilizer_l_symbol}) of the
              $2$-cocycle $\psi(a,b)=R^{x_0}_{a,b}$, with
              $\gamma^x_g=R^{x_0}_{k_{gx},t_k(g,x)}/R^{x_0}_{g,k_x}$.
        \item If $\omega$ is a $3$-cocycle equal to one on
              $H\times H\times H$, every $L$-symbol compatible with $\omega$
              is action-gauge equivalent to
              (\ref{eq:mpug_stabilizer_reconstruction}) for some $2$-cocycle
              $\psi$ on $H$. For a general scalar $2$-cochain $\beta$ making
              $\beta\mathbin{\triangleright}\omega$ trivial on $H$, the
              same holds for $(\beta,1)\mathbin{\triangleright}L$ and the
              formula built from $\beta\mathbin{\triangleright}\omega$.
        \item Let $L_0$ be compatible with an arbitrary scalar $3$-cochain
              $\omega$. An $L$-symbol $L$ is compatible with $\omega$ if and
              only if $L$ is action-gauge equivalent to $L_0L[\psi]$ for
              some $2$-cocycle $\psi$ on $H$, and $L_0L[\psi_1]$ and
              $L_0L[\psi_2]$ are action-gauge equivalent if and only if
              $\psi_1\sim\psi_2$. Hence sending the class of $\psi$ to the
              class of $L_0L[\psi]$ is a bijection from $H^2(H,\C^\times)$
              onto the action-gauge classes of $L$-symbols compatible with
              $\omega$.
        \item Sending the class of $\psi$ and the class of $L$ to the
              class of $L\,L[\psi]$ is a well-defined action of the group
              cohomology $H^2(H,\C^\times)$, for the trivial action of $H$
              on $\C^\times$, on these classes, and for any two classes
              exactly one cohomology class carries the first to the second.
    \end{enumerate}
    The last two items are the torsor statement of \cite{GarreRubio2023MPOSym}
    for fixed $(G,\omega,H)$, with the two-cochain held fixed as in the
    source. In all four items $G$, and hence $H$, is a small group: its
    underlying set lies in a fixed set-theoretic universe containing $\Z$,
    in which the group cohomology $H^2(H,\C^\times)$ is formed. In the last
    two items the set $\mathsf X$ lies in the same universe.
    % Source anchor: arXiv:2203.12563, REsubmission.tex lines 749--761.
\end{lemma}
\begin{proof}\leanok
    \uses{thm:mpug_stabilizer_reconstruction, thm:mpug_stabilizer_cocycle_l_symbol}
    For the first item, compatibility at $x_0$ with arguments
    $(g,h,k_x)$ expresses $R^x_{g,h}$ through basepoint values. Rewriting
    $hk_x=k_{hx}t_k(h,x)$ and $gk_{hx}=k_{ghx}t_k(g,hx)$ and using
    compatibility at $x_0$ with arguments $(g,k_{hx},t_k(h,x))$ and
    $(k_{ghx},t_k(g,hx),t_k(h,x))$ combines the three equations into
    $R^x_{g,h}=\gamma^x_{gh}(\gamma^{hx}_g\gamma^x_h)^{-1}
        \psi(t_k(g,hx),t_k(h,x))$. Compatibility at $x_0$ with arguments in
    $H$ is the cocycle equation for $\psi$.

    For the third item, if $L$ and $L_0$ are compatible with $\omega$,
    then $LL_0^{-1}$ is compatible with the trivial $3$-cochain, and the first
    item applies; conversely $L_0L[\psi]$ is compatible with $\omega$, and an
    gauge with trivial two-cochain preserves compatibility and
    commutes with multiplication by $L_0$. Cancelling $L_0$ reduces the
    freeness statement to (\ref{eq:mpug_induced_cohomology_iff}). The second
    item is the third with $L_0$ the formula at $\psi=1$, which is
    compatible by Lemma~\ref{thm:mpug_stabilizer_reconstruction}.
    For the last item, cocycles modulo cohomology form the second group
    cohomology of $H$ with coefficients in $\C^\times$, and
    \begin{align}
        [\psi_1]+[\psi_2] & =[\psi_1\psi_2],
        \notag                                     \\
        L[\psi_1\psi_2]   & =L[\psi_1]\,L[\psi_2].
        \notag
    \end{align}
    Cohomologous cocycles give induced $L$-symbols related by an action
    gauge, so the rule is well defined, and by the two identities acting by
    $[\psi_1]+[\psi_2]$ is acting by $[\psi_2]$ and then by $[\psi_1]$. With
    base point $L$ the action is the bijection of the third item.
\end{proof}

\begin{lemma}[Explicit gauge from a stabilizer cocycle extension]
    \label{thm:mpug_stabilizer_extension_gauge}
    \lean{StabilizerRepresentatives.extensionActionGauge,
        StabilizerRepresentatives.inducedLSymbol_eq_extensionActionGauge,
        StabilizerRepresentatives.inducedLSymbol_hasTrivialGauge_of_extension}
    \leanok
    \uses{thm:mpug_stabilizer_cocycle_l_symbol, def:mpug_l_gauge, def:mpug_block_independence_gauge_forms}
    Let $G$ act on $\mathsf X$, with basepoint $x_0$, stabilizer
    $H=\operatorname{Stab}_G(x_0)$, and representatives $k_x$ satisfying
    $k_xx_0=x$ and $k_{x_0}=e$.
    Let $\psi\colon H\times H\to\C^\times$ and let
    $\Psi\colon G\times G\to\C^\times$ be a two-cocycle whose restriction
    to $H\times H$ is $\psi$. Define
    \begin{align}
        \gamma_g^x
         & =\frac{\Psi(k_{gx}^{-1},gk_x)\Psi(g,k_x)}
            {\Psi(e,e)\Psi(k_x,k_x^{-1})}.
        \label{eq:mpug_extension_gamma}
    \end{align}
    With $h_1=t_k(g_1,g_2x)$ and $h_2=t_k(g_2,x)$, the induced symbols satisfy
    \begin{align}
        L^x_{g_1,g_2}=\psi(h_1,h_2)
         & =\Psi(g_1,g_2)\frac{\gamma_{g_1g_2}^x}
                         {\gamma_{g_1}^{g_2x}\gamma_{g_2}^x}.
        \label{eq:mpug_extension_identity}
    \end{align}
    Consequently the inverse gauges
    $(\Psi^{-1},\gamma^{-1})$ make $L$ identically one.
    This is the scalar representative calculation in the forward proof of
    the block-independence extension proposition of
    \cite{FrancoRubio2025MPUGauging};
    here extension means the equality $\Psi|_{H\times H}=\psi$ of chosen
    cocycle representatives, rather than only equality of their classes.
\end{lemma}
\begin{proof}\leanok
    With $a=k_{g_1g_2x}$, $b=k_{g_2x}$, $c=k_x$, $u=g_1$, and $v=g_2$,
    the cocycle equation gives $\Psi(z,e)=\Psi(e,z)=\Psi(e,e)$ and
    $\Psi(b^{-1},b)=\Psi(b,b^{-1})$. Applying it successively to the triples
    \begin{align}
         & (a^{-1}ub,b^{-1},vc),\quad (a^{-1}ub,b^{-1},b),\quad
        (a^{-1},u,vc),\quad (a^{-1},u,b),\quad (u,v,c)
    \end{align}
    and solving these equalities for the corresponding factors of
    $\Psi(a^{-1}ub,b^{-1}vc)$ yields
    (\ref{eq:mpug_extension_identity}) with
    (\ref{eq:mpug_extension_gamma}) after cancellation. The restriction
    hypothesis identifies its left side with $\psi(h_1,h_2)$. Substitution into
    (\ref{eq:mpug_l_gauge}) cancels the multiplier under the inverse gauges.
\end{proof}

\begin{lemma}[Exact extension from trivializing scalar gauges]
    \label{lem:mpug_trivializing_extension}
    \lean{LSymbol.isMulCocycle₂_of_isTrivial_gauge,
        LSymbol.basepointExtension,
        LSymbol.basepointExtension_isMulCocycle₂,
        StabilizerRepresentatives.basepointExtension_restrict,
        StabilizerRepresentatives.exists_extension_of_inducedLSymbol_hasTrivialGauge}
    \leanok
    \uses{thm:mpug_stabilizer_cocycle_l_symbol, def:mpug_l_gauge}
    Let $G$ act on $\mathsf X$ with basepoint $x_0$, stabilizer
    $H=\operatorname{Stab}_G(x_0)$, and representatives $k_x$ satisfying
    $k_xx_0=x$ and $k_{x_0}=e$. Let $\psi:H\times H\to\C^\times$ be a
    two-cocycle and $L[\psi]$ its induced scalar $L$-symbol.
    Suppose scalar cochains $\beta:G\times G\to\C^\times$ and
    $\gamma:G\times\mathsf X\to\C^\times$ satisfy
    $(\beta,\gamma)\mathbin{\triangleright}L[\psi]=1$.
    Then $\beta$ is a two-cocycle and
    \begin{align}
        \Psi(g,h)
         & =\beta(g,h)^{-1}
        \frac{\gamma_{g,x_0}\gamma_{h,x_0}}{\gamma_{gh,x_0}}
        \label{eq:mpug_trivializing_extension}
    \end{align}
    is a two-cocycle with $\Psi|_{H\times H}=\psi$ exactly.
    Neither cocycles nor gauges need be normalized.
    This is the converse scalar calculation in
    \cite[lines~7057--7064]{FrancoRubio2025MPUGauging}, with gauges that
    trivialize $L[\psi]$, the inverses of the source's forward gauges.
\end{lemma}
\begin{proof}\leanok
    \uses{thm:mpug_stabilizer_cocycle_l_symbol, def:mpug_l_gauge, thm:mpug_h2_exact_extension}
    Compatibility of $L[\psi]$ with $\omega=1$ transports under
    $(\beta,\gamma)$ to compatibility of the constant symbol $1$ with
    $d\beta$. Thus $d\beta=1$, which is the two-cocycle equation for
    $\beta$. Setting $\varphi(g)=\gamma_{g,x_0}$, formula
    (\ref{eq:mpug_trivializing_extension}) is
    $\Psi=d\varphi\,\beta^{-1}$, hence is a two-cocycle.
    For $a,b\in H$, the equality $bx_0=x_0$ and the chosen
    $k_{x_0}=e$ identify $L[\psi]^{x_0}_{a,b}$ with $\psi(a,b)$.
    Evaluating the trivializing-gauge identity there gives
    $\Psi(a,b)=\psi(a,b)$.
\end{proof}

\begin{lemma}[Block independence extends the stabilizer cocycle]
    \label{lem:mpug_bi_basepoint_extension}
    \lean{LSymbol.basepointExtension_apply_of_isTrivial_gauge,
        LSymbol.exists_extension_of_isBlockIndependent}
    \leanok
    \uses{lem:mpug_trivializing_extension, thm:mpug_block_independence_four_forms, def:mpug_block_independence_gauge_forms}
    Let $G$ act on $\mathsf X$ with basepoint $x_0$ and stabilizer
    $H=\operatorname{Stab}_G(x_0)$, and let $L$ be any $L$-symbol compatible
    with $\omega=1$. If $L$ is block independent after a joint scalar gauge,
    then there is a two-cocycle $\Psi:G\times G\to\C^\times$ with
    $\Psi(a,b)=L^{x_0}_{a,b}$ for all $a,b\in H$. This is the converse
    direction of \cite[lines~7057--7064]{FrancoRubio2025MPUGauging} for an
    arbitrary compatible $L$-symbol, not only for one induced by
    representatives.
\end{lemma}
\begin{proof}\leanok
    \uses{lem:mpug_trivializing_extension, thm:mpug_block_independence_four_forms}
    By Lemma~\ref{thm:mpug_block_independence_four_forms}, some joint scalar
    gauge $(\beta,\gamma)$ trivializes $L$. The cochain
    (\ref{eq:mpug_trivializing_extension}) is a two-cocycle by
    Lemma~\ref{lem:mpug_trivializing_extension}. For $b\in H$ the trivializing
    identity at $(x_0,a,b)$ uses only $bx_0=x_0$, and it gives
    $\Psi(a,b)=L^{x_0}_{a,b}$.
\end{proof}

\begin{lemma}[Restriction is the inhomogeneous restriction map]
    \label{lem:mpug_restriction_inhomogeneous}
    \lean{ScalarCocycle.toInhomogeneousCochain_comap,
        ScalarThreeCochain.toInhomogeneousCochain_comap}
    \leanok
    \uses{def:mpug_h2_pullback, def:mpug_cochain_restriction, def:mpug_h3_group_cohomology}
    Let $f:K\to G$ be a group homomorphism. Under the identification of scalar
    $2$- and $3$-cochains with inhomogeneous cochains valued in the trivial
    representation $\C^\times_{\mathrm{triv}}$, the pullback $f^*$ in degrees
    two and three is the restriction map of inhomogeneous cochains along $f$,
    $x\mapsto x\circ(f\times\cdots\times f)$, whose coefficient map is the
    identity of $\C^\times$.
\end{lemma}
\begin{proof}\leanok
    \uses{def:mpug_h2_pullback, def:mpug_cochain_restriction, def:mpug_h3_group_cohomology}
    Both sides send a cochain $x$ to $(g_1,\dots,g_n)\mapsto x(f(g_1),\dots,f(g_n))$.
\end{proof}

\begin{lemma}[Scalar block independence and restriction of the stabilizer class]
    \label{lem:mpug_scalar_bi_restriction}
    \lean{StabilizerRepresentatives.inducedLSymbol_isBlockIndependent_iff_exists_extension,
        StabilizerRepresentatives.inducedLSymbol_isBlockIndependent_iff_H2π_mem_range_map_subtype}
    \leanok
    \uses{lem:mpug_trivializing_extension, thm:mpug_stabilizer_extension_gauge, thm:mpug_h2_exact_extension, thm:mpug_block_independence_four_forms}
    Let $G$ act on $\mathsf X$ with basepoint $x_0$, stabilizer
    $H=\operatorname{Stab}_G(x_0)$, and representatives $k_x$ satisfying
    $k_xx_0=x$ and $k_{x_0}=e$. For a two-cocycle
    $\psi:H\times H\to\C^\times$, the following are equivalent:
    \begin{enumerate}
        \item The induced scalar $L$-symbol $L[\psi]$ is block independent
              after a joint scalar gauge.
        \item There is a two-cocycle $\Psi:G\times G\to\C^\times$ with
              $\Psi|_{H\times H}=\psi$ exactly.
        \item The class $[\psi]$ belongs to the image of the restriction map
              $H^2(G,\C^\times)\to H^2(H,\C^\times)$ for trivial coefficients.
    \end{enumerate}
    No normality or finiteness assumption is imposed on $H$ or $G$.
    This is the $L$-symbol form of
    \cite[lines~7042--7064]{FrancoRubio2025MPUGauging} for the induced symbols.
\end{lemma}
\begin{proof}\leanok
    \uses{thm:mpug_block_independence_four_forms, lem:mpug_trivializing_extension, thm:mpug_stabilizer_extension_gauge, thm:mpug_h2_exact_extension}
    By Lemma~\ref{thm:mpug_block_independence_four_forms}, scalar block
    independence is equivalent to the existence of a trivializing scalar
    gauge. Lemma~\ref{lem:mpug_trivializing_extension} constructs an exact
    extension from such a gauge, and
    Lemma~\ref{thm:mpug_stabilizer_extension_gauge} constructs a trivializing
    gauge from an exact extension. Finally,
    Lemma~\ref{thm:mpug_h2_exact_extension} identifies exact extension of
    $\psi$ with membership of $[\psi]$ in the image of restriction.
\end{proof}

\begin{lemma}[Scalar block independence with circle coefficients]
    \label{lem:mpug_scalar_bi_restriction_circle}
    \lean{ScalarCohomology.circleRepresentationRes,
        ScalarCohomology.circleToComplexUnits_map_comp_map_res,
        ScalarCohomology.circleToComplexUnits_map_res_apply,
        ScalarCohomology.mem_range_map_circleRepresentationRes_iff,
        LSymbol.isBlockIndependent_gauge_iff,
        StabilizerRepresentatives.inducedLSymbol_isBlockIndependent_iff_circle}
    \leanok
    \uses{lem:mpug_scalar_bi_restriction, thm:mpug_circle_complex_units_h2}
    Restriction along a group homomorphism $f\colon S\to G$ commutes with
    the maps $H^2(-,U(1))\to H^2(-,\C^\times)$ induced by inclusion of the
    unit circle. If $G$ is finite, a class $c\in H^2(S,U(1))$ is the
    restriction of a class of $G$ exactly when its image in
    $H^2(S,\C^\times)$ is. Block independence of an $L$-symbol after
    a joint scalar gauge is unchanged by any further joint scalar gauge.

    Consequently, in the setting of Lemma~\ref{lem:mpug_scalar_bi_restriction}
    with $G$ finite, the induced $L$-symbol $L[\psi]$ is block independent
    after a joint scalar gauge if and only if the class in $H^2(H,U(1))$ of
    the pointwise phase $\psi/|\psi|$ lies in the image of the restriction map
    $H^2(G,U(1))\to H^2(H,U(1))$. For a cocycle of modulus one this phase is
    $\psi$ itself. This is the coefficient group of
    \cite[lines~7042--7064]{FrancoRubio2025MPUGauging}; the comparison of
    coefficients is the algebraic one of
    Lemma~\ref{thm:mpug_circle_complex_units_h2}.
\end{lemma}
\begin{proof}\leanok
    \uses{lem:mpug_scalar_bi_restriction, thm:mpug_circle_complex_units_h2, thm:mpug_l_gauge_composition}
    Both composites of restriction and inclusion are induced by the same
    homomorphism $S\to G$ and the same coefficient map
    $U(1)\to\C^\times$, so they agree. If $c=\operatorname{res}(a)$, then
    the image of $c$ is the restriction of the image of $a$. Conversely, if
    the image of $c$ is $\operatorname{res}(b)$, surjectivity of inclusion
    for the finite group $G$ writes $b$ as the image of some $a$; then $c$
    and $\operatorname{res}(a)$ have the same image, and inclusion is
    injective for every group. Gauge invariance follows from
    (\ref{eq:mpug_l_gauge_composition}) applied with the inverse gauge. By
    Lemma~\ref{thm:mpug_circle_complex_units_h2} the class of the phase of
    $\psi$ maps to $[\psi]$, and Lemma~\ref{lem:mpug_scalar_bi_restriction}
    finishes the proof.
\end{proof}
